\subsection{DEFINITION OF THE ENVELOPING ALGEBRA}

Let g be a Lie algebra over K. For any associative algebra with unit element L over K, an \(\alpha\) -mapping of g into L is a K-linear mapping \(\sigma\) of g into L such that

\[
\sigma ([ x, y ]) = \sigma (x) \sigma (y) - \sigma (y) \sigma (x) \quad (x, y \text {   in   } g)
\]

(in other words a homomorphism of \(\mathfrak{g}\) into the Lie algebra associated with L). If L' is another associative algebra with unit element over K and \(\tau\) a homomorphism of L into L' mapping 1 to 1, then \(\tau \circ \sigma\) is an \(\alpha\) -mapping of \(\mathfrak{g}\) into L'. We shall look for an associative algebra with unit element and an \(\alpha\) -mapping of \(\mathfrak{g}\) into this algebra which are universal (Set Theory, Chapter IV, § 3, no. 1).

DEFINITION 1. Let \(\mathfrak{g}\) be a Lie algebra over \(\mathbf{K}\) , \(\mathbf{T}\) the tensor algebra of the \(\mathbf{K}\) -module \(\mathfrak{g}\) and \(\mathbf{J}\) the two-sided ideal of \(\mathbf{T}\) generated by the tensors \(x \otimes y - y \otimes x - [x, y]\) where \(x \in \mathfrak{g}, y \in \mathfrak{g}\) . The associative algebra \(\mathbf{U} = \mathbf{T} / \mathbf{J}\) is called the enveloping algebra of \(\mathfrak{g}\) . The restriction to \(\mathfrak{g}\) of the canonical mapping of \(\mathbf{T}\) onto \(\mathbf{U}\) is called the canonical mapping of \(\mathfrak{g}\) into \(\mathbf{U}\) .

Let \(T_{+}\) be the two-sided ideal of T consisting of the tensors whose component of order 0 is zero. Let \(\mathrm{T}_0 = \mathrm{K}.1\) be the set of elements of T of order 0. Let \(\mathrm{U}_{+}\) and \(\mathrm{U}_0\) be the canonical images of \(\mathrm{T}_{+}\) and \(\mathrm{T}_0\) in U. As \(\mathrm{J} \subset \mathrm{T}_{+}\) , the decomposition into a direct sum \(\mathrm{T} = \mathrm{T}_0 + \mathrm{T}_{+}\) implies a decomposition into a direct sum \(\mathrm{U} = \mathrm{U}_0 + \mathrm{U}_{+}\) . The algebra U therefore has a unit element distinct from 0 and \(\mathrm{U}_0 = \mathrm{K}.1\) . For all \(x \in \mathbf{U}\) , the component of \(x\) in \(\mathrm{U}_0\) is called the constant term of \(x\) . The elements with constant term zero form a two-sided ideal of U, namely the two-sided ideal \(\mathrm{U}^{+}\) generated by the canonical image of g in U.

The associative algebra U is generated by 1 and the canonical image of g in U.

If \(x \in \mathfrak{g}\) and \(y \in \mathfrak{g}\) , \(x \otimes y - y \otimes x\) and \([x, y]\) are congruent in T modulo J; hence, if \(\sigma_0\) denotes the canonical mapping of \(\mathfrak{g}\) into U,

\[
\sigma_ {0} (x) \sigma_ {0} (y) - \sigma_ {0} (y) \sigma_ {0} (x) = \sigma_ {0} ([ x, y ])
\]

in U. In other words, \(\sigma_{0}\) is an \(\alpha\) -mapping of g into U.

PROPOSITION 1. Let \(\sigma\) be an \(\alpha\) -mapping of \(\mathfrak{g}\) into the associative algebra \(\mathbf{L}\) with unit element. There exists one and only one homomorphism \(\tau\) of \(\mathbf{U}\) into \(\mathbf{L}\) , mapping 1 to 1, such that \(\sigma = \tau \circ \sigma_0\) , where \(\sigma_0\) denotes the canonical mapping of \(\mathfrak{g}\) into \(\mathbf{U}\) .

Let \(\tau'\) be the unique homomorphism of T into L which extends \(\sigma\) and maps 1 to 1. Then, for \(x, y\) in g,

\[
\tau^ {\prime} (x \otimes y - y \otimes x - [ x, y ]) = \sigma (x) \sigma (y) - \sigma (y) \sigma (x) - \sigma ([ x, y ]) = 0;
\]

hence \(\tau'\) is zero on \(J\) and defines on passing to the quotient a homomorphism \(\tau\) of \(U\) into \(L\) , mapping 1 to 1, such that \(\sigma = \tau \circ \sigma_0\) . The uniqueness of \(\tau\) is immediate since \(\sigma_0(g)\) and 1 generate the algebra \(U\) .

Let \(g'\) be another Lie algebra over K, U' its enveloping algebra and \(\sigma_0'\) the canonical mapping of \(g'\) into U'. Let \(\phi\) be a homomorphism of g into g'. Then \(\sigma_0' \circ \phi\) is an \(\alpha\) -mapping of g into U'; hence there exists one and only one homomorphism \(\tilde{\phi}\) of U into U' mapping 1 to 1 and such that the diagram

\[
\begin{array}{c} g \xrightarrow {\phi} g ^ {\prime} \\ \sigma_ {0} \Biggl \downarrow \\ U \xrightarrow {\tilde {\phi}} U ^ {\prime} \end{array}
\]

is commutative. This homomorphism maps the elements of U whose constant term is zero to elements of U' whose constant term is zero. If g'' is another Lie algebra over K and \(\phi'\) is a homomorphism of g' into g'', then \(\widetilde{\phi'\circ\phi}=\widetilde{\phi}'\circ\widetilde{\phi}\) .

